% 量子力学1 量子力学1-11 問題7（086-symwell） 解いた日 2026-11-20
%% kotae
\toi{1}{○}{記述}
\begin{align*}
&\phi\left(\frac{a}{2}\right)=A\cos\frac{ka}{2}+B\sin\frac{ka}{2}=0\\
&\phi\left(-\frac{a}{2}\right)=A\cos\frac{ka}{2}-B\sin\frac{ka}{2}=0\\
&\therefore\ 2A\cos\frac{ka}{2}=0,\quad B\sin\frac{ka}{2}=0
\end{align*}
\begin{align*}
&\cos\frac{ka}{2}=0\qquad\therefore\ \frac{ka}{2}=\frac{\pi}{2}+n\pi\\
&\therefore\ \frac{\sqrt{2mE_{n}}}{\hbar}=\frac{1}{a}(2n+1)\pi\\
&E_{n}=\frac{(2n+1)^{2}\pi^{2}\hbar^{2}}{2ma^{2}}
\end{align*}
\begin{align*}
&\sin\frac{ka}{2}=0\qquad\therefore\ \frac{ka}{2}=n\pi\\
&\frac{\sqrt{2mE_{n}}}{\hbar}=\frac{2n\pi}{a}\\
&E_{n}=\frac{(2n)^{2}\pi^{2}\hbar^{2}}{2ma^{2}}
\end{align*}
\[ \therefore\ E_{n}=\frac{n^{2}\pi^{2}\hbar^{2}}{2ma^{2}}\quad\nuke{n=1,2,3,\dots} \]
\naze{最後の $n$ を正の整数に限る理由が無い。$k>0$ で，$A=B=0$（$\phi=0$）を除くので $n=1,2,3,\dots$ である。なお，2つの場合で同じ文字 $n$ を別の意味に使っている（上の2式の $2n+1$ と $2n$ が最後の式の $n$ にあたる）。}
$n$ が奇数のとき、
\[ \cos\frac{ka}{2}=0\qquad B=0 \]
\[ \chi_{n}(x)=A\cos\frac{n\pi x}{a} \]
\begin{align*}
\int_{-\infty}^{\infty}\abs{\chi_{n}(x)}^{2}\dd x&=A^{2}\int_{-\frac{a}{2}}^{\frac{a}{2}}\cos^{2}\frac{n\pi x}{a}\dd x\\
&=\frac{A^{2}}{2}\int_{-\frac{a}{2}}^{\frac{a}{2}}1+\cos\frac{2n\pi x}{a}\dd x\\
&=\frac{A^{2}}{2}\left[x+\frac{a}{2n\pi}\sin\frac{2n\pi x}{a}\right]_{-\frac{a}{2}}^{\frac{a}{2}}\\
&=\frac{A^{2}}{2}\left(\frac{a}{2}-\left(-\frac{a}{2}\right)\right)
\end{align*}
\[ \frac{aA^{2}}{2}=1\qquad A=\sqrt{\frac{2}{a}} \]
\[ \chi_{n}(x)=\sqrt{\frac{2}{a}}\cos\frac{n\pi x}{a}\quad(\text{$n$ が奇数}) \]
$n$ が偶数のとき、
\[ \sin\frac{ka}{2}=0\qquad A=0 \]
\[ \chi_{n}(x)=\sqrt{\frac{2}{a}}\sin\frac{n\pi x}{a}\quad(\text{$n$ は偶数}) \]
$n$ が奇数のとき偶パリティ

$n$ が偶数のとき奇パリティ
\toi{2}{◎}{}
\begin{align*}
\phi_{n}\left(x+\frac{a}{2}\right)&=\sqrt{\frac{2}{a}}\sin\frac{n\pi\left(x+\frac{a}{2}\right)}{a}\\
&=\sqrt{\frac{2}{a}}\sin\frac{n\pi x}{a}+\frac{n\pi}{2}
\end{align*}
$n$ が奇数のとき、
\[ \phi_{n}\left(x+\frac{a}{2}\right)=\pm\sqrt{\frac{2}{a}}\cos\frac{n\pi x}{a} \]
$n$ が偶数のとき、
\[ \phi_{n}\left(x+\frac{a}{2}\right)=\pm\sqrt{\frac{2}{a}}\sin\frac{n\pi x}{a} \]
符号を除いて一致
\toi{3}{△}{計算}
i) $n$ が奇数のとき
\begin{align*}
\expval{x}&=\int_{-\frac{a}{2}}^{\frac{a}{2}}\underbrace{x\cdot\frac{2}{a}\cos^{2}\frac{n\pi x}{a}}_{\text{奇関数}}\dd x\\
&=0
\end{align*}
\begin{align*}
\expval{x^{2}}&=\int_{-\frac{a}{2}}^{\frac{a}{2}}x^{2}\cdot\frac{2}{a}\cos^{2}\frac{n\pi x}{a}\dd x\\
&=\frac{4}{a}\int_{0}^{\frac{a}{2}}x^{2}\cos^{2}\frac{n\pi x}{a}\dd x\\
&=\frac{2}{a}\int_{0}^{\frac{a}{2}}x^{2}+x^{2}\cos\frac{2n\pi x}{a}\dd x\\
&=\frac{2}{a}\cdot\frac{1}{3}\left(\frac{a}{2}\right)^{3}+\frac{2}{a}\int_{0}^{\frac{a}{2}}x^{2}\cos\frac{2n\pi x}{a}\dd x
\end{align*}
（矢印で第2項の積分につなぐ）
\begin{align*}
&\left[x^{2}\cdot\frac{a}{2n\pi}\sin\frac{2n\pi x}{a}\right]_{0}^{\frac{a}{2}}\\
&\qquad-\int_{0}^{\frac{a}{2}}\ayamari{x}{2x}\cdot\frac{a}{2n\pi}\sin\frac{2n\pi x}{a}\dd x\\
&=-\frac{a}{2n\pi}\int_{0}^{\frac{a}{2}}x\sin\frac{2n\pi x}{a}\dd x
\end{align*}
\naze{部分積分 $\int uv'\dd x=[uv]-\int u'v\dd x$ で $u=x^{2}$ の微分は $2x$ である。正しくは第2項が $-(a/(n\pi))\int_{0}^{a/2}x\sin(2n\pi x/a)\dd x$ で，以降はこの誤り（第2項が半分）を引き継ぐ。正しい答えは $\expval{x^{2}}=a^{2}/12-a^{2}/(2n^{2}\pi^{2})$（$n$ が偶数でも同じ）である。}
\begin{align*}
&\int_{0}^{\frac{a}{2}}x\sin\frac{2n\pi x}{a}\dd x\\
&\quad=\left[-\frac{ax}{2n\pi}\cos\frac{2n\pi x}{a}\right]_{0}^{\frac{a}{2}}+\int_{0}^{\frac{a}{2}}\frac{a}{2n\pi}\cos\frac{2n\pi x}{a}\dd x\\
&\quad=-\frac{a^{2}}{4n\pi}(-1)^{n}+\left(\frac{a}{2n\pi}\right)^{2}\left[\sin\frac{2n\pi x}{a}\right]_{0}^{\frac{a}{2}}
\end{align*}
\begin{align*}
\expval{x^{2}}&=\frac{a^{2}}{12}+\frac{2}{a}\left(-\frac{a}{2n\pi}\cdot-\frac{a^{2}}{4n\pi}(-1)^{n}\right)\\
&=\frac{a^{2}}{12}+\frac{a^{2}}{4n^{2}\pi^{2}}(-1)^{n}
\end{align*}
$n$ は奇数なので
\[ \expval{x^{2}}=\frac{a^{2}}{12}-\frac{a^{2}}{4n^{2}\pi^{2}} \]
ii) $n$ が偶数のとき、
\begin{align*}
\expval{x}&=\int_{-\frac{a}{2}}^{\frac{a}{2}}\underbrace{x\cdot\frac{2}{a}\sin^{2}\frac{n\pi x}{a}}_{\text{奇関数}}\dd x=0
\end{align*}
\begin{align*}
\expval{x^{2}}&=\int_{-\frac{a}{2}}^{\frac{a}{2}}x^{2}\cdot\frac{2}{a}\sin^{2}\frac{n\pi x}{a}\dd x\\
&=\frac{4}{a}\int_{0}^{\frac{a}{2}}x^{2}\sin^{2}\frac{n\pi x}{a}\dd x\\
&=\ayamari{\frac{4}{a}}{\frac{2}{a}}\int_{0}^{\frac{a}{2}}x^{2}-x^{2}\cos\frac{2n\pi x}{a}\dd x\\
&=\frac{a^{2}}{12}-\frac{a^{2}}{4n^{2}\pi^{2}}(-1)^{n}\\
&=\frac{a^{2}}{12}-\frac{a^{2}}{4n^{2}\pi^{2}}
\end{align*}
\naze{半角の公式 $\sin^{2}\theta=(1-\cos2\theta)/2$ の $1/2$ で係数は $2/a$ になる（次の行の $a^{2}/12$ は $2/a$ で計算していて，答えには影響しない）。第2項は i) と同じ部分積分の誤りを引き継ぎ，正しくは $-a^{2}/(2n^{2}\pi^{2})$ である。}
\[ \expval{x}=0\qquad\expval{x^{2}}=\frac{a^{2}}{12}-\frac{a^{2}}{4n^{2}\pi^{2}} \]
\[ n\to\infty\ \text{で}\ \expval{x^{2}}\to\frac{a^{2}}{12} \]
一様に分布しているとき、規格化条件より、
\[ \phi(x)=\frac{1}{a} \]
\begin{align*}
\int_{-\frac{a}{2}}^{\frac{a}{2}}x^{2}\phi(x)\dd x&=\frac{2}{a}\int_{0}^{\frac{a}{2}}x^{2}\dd x\\
&=\frac{2}{a}\left[\frac{1}{3}x^{3}\right]_{0}^{\frac{a}{2}}\\
&=\frac{a^{2}}{12}\quad\text{となり一致する。}
\end{align*}
%% kizuki
\item 部分積分 $\int uv'\dd x=[uv]-\int u'v\dd x$ は，$u=x^{2}$，$v'=\cos qx$ と役割を書いてから $u'=2x$ を代入する。2回続けるときは1回ごとに係数を書き出す。
\item 検算は，得た原始関数 $\frac{x^{2}}{q}\sin qx+\frac{2x}{q^{2}}\cos qx-\frac{2}{q^{3}}\sin qx$ を微分して $x^{2}\cos qx$ に戻るかで行う。$n\to\infty$ の極限は第2項の係数によらないので，係数の誤りの確かめにならない。
\item \mae{第4章 問題3・10-28} (1) の $n$ の範囲は問題4 (1) と同じ。2つの場合を別の文字（$n=2l+1$，$n=2l$ など）で書くと，最後にまとめるときに範囲を書き落としにくい。
